In driving and other epoch-structured paradigms, a common approach is to apply a continuous-signal blink detector to the raw or minimally preprocessed recording and then map the detected blinks back onto the epoch grid. We refer to this as the pooled, continuous-signal strategy, because each channel's detection threshold is estimated once from a signal pool that contains both blink-heavy and blink-free epochs. BLINKER~\citep{kleifges2017blinker} and the EOG event detector of MNE-Python~\citep{GramfortEtAl2013a} both follow this strategy, although they estimate the threshold in different ways.
BLINKER combines the mean of the pooled signal with a robust dispersion term, so blink-free epochs dilute these amplitude statistics and pull the threshold down. The MNE detector instead sets its threshold to one quarter of the range of the filtered signal, so a single large excursion anywhere in the recording raises the threshold for every epoch. The premise of the proposed pipeline is that either form of pooling misplaces the threshold for part of the recording and thereby shifts the resulting precision-recall trade-off.
